\section{Task suite}
\label{sec:tasks}

\NTasks{} tasks cover \NTaskAlgorithms{} algorithms in \NFamilies{} families, two variants
each (Appendix Table~\ref{tab:tasks}). A task states a rule in prose; the model implements
\texttt{solve(prices, split\_factors)} returning positions, or for volatility, risk and
forecast tasks the value known at the previous close. Each task can be solved by hand or by
one library call. The rules are the momentum, volatility-scaling and allocation baselines
that learned strategies are measured against
\citep{lim2019enhancing,wood2022slow,zhang2020deep}.

\paragraph{Data.}
Synthetic panels have four assets and \NRows{} business days of seeded Gaussian returns,
with four splits at fixed rows (2-for-1, 3-for-1, 3-for-2, 1-for-2), on \NSeeds{} seeds.
Variant~a gives cumulative factors in alphabetical order; variant~b gives event ratios,
shuffled columns and other parameters. The suite also has a real panel: as-traded
closes of Amazon, Shopify, Alphabet and Tesla for \NRealRows{} sessions from \RealFirst{} to
\RealLast{}, in which all four split (20-, 10-, 20- and 3-for-1; Appendix
Figure~\ref{fig:real}). Yahoo Finance serves the
closes split-adjusted; we rebuild as-traded quotes from a split table checked against the
companies' announcements and published pre-split closes, and release the fetch code and the
file's SHA-256, since the prices may not be redistributed.

\paragraph{Judge and hidden grader.}
The judge runs each program in a separate process with a memory limit and a 120-second
timeout (in our runs also confined with Landlock and seccomp), refuses \texttt{import fin\_skills} without library access, and applies the three
checks of \S\ref{sec:library}. A program is \emph{valid} only if it
ran, passed every check and matches an independent numpy reference within $10^{-8}$
($10^{-5}$ for \texttt{min\_variance}). A program that passes the checks but is wrong is
\emph{accepted but wrong}; we report that rate because the checks cannot see it.

\paragraph{Qualification.}
Before any model ran: on every task and data panel, the reference agrees with the library
to within $\RefGapOther$ ($\RefGapMinVariance$ for \texttt{min\_variance}); a correct
program passes every check and is graded valid; and on a position task and a table task,
planted versions of each failure (same close, no adjustment, factors cumulated twice,
missing warm-up) fail their check and the grade. Building the references exposed a library
defect: the long-only minimum-variance optimiser stopped early on daily-scale covariances,
off by up to $1.2\times10^{-2}$, until we rescaled the covariance before solving.

\paragraph{Split.}
The split is by algorithm. The \NForcedDev{} algorithms a model has already seen go to
development, and a seeded draw sends \NDrawnDev{} of the other \NDrawn{} there too:
\NDevTasks{} development and \NHeldTasks{} held-out tasks, the latter covering all
\NFamilies{} families.

\section{Study protocol}
\label{sec:study}

We fixed the tasks, prompts, seeds, budgets, models and contrasts before any test run, on
the suite as it stood on 8~November 2026: \StudyDevTasks{} development and
\StudyHeldTasks{} held-out position tasks on synthetic panels (Bollinger, RSI and Donchian
rules, two variants each). The tasks added since, and the real panel, are not part of the
study. It has three tracks.

\paragraph{P: feedback and library.}
The three controlled arms, with and without the library, on the \StudyHeldTasks{} held-out
tasks and \StudySeeds{} seeds: \StudyPEpisodes{} episodes per model. Three contrasts were
declared: \texttt{Check} against \texttt{Repeat} and \texttt{Explain} against
\texttt{Check}, both without the library, and library against none under \texttt{Explain}.

\paragraph{E: external questions.}
\StudyEQuestions{} public questions from the two table-and-text program-synthesis tasks of
BizBench \citep{krumdick2024bizbench} (\StudyECodeFinQA{} CodeFinQA, \StudyECodeTATQA{}
CodeTAT-QA), selected by a rule fixed before any model output: a numeric answer, a prompt
that fits every core model's context, no duplicate in the training split, one question per
group of shared context. The model writes one program per question, with and without the
library and without feedback: \StudyEEpisodes{} episodes per model. An answer counts if the
program runs and returns the public label exactly; BizBench scores within 1\%, which we report as
a sensitivity analysis.

\paragraph{S: seven methods.}
The \NMethods{} methods of \S\ref{sec:methods}, with and without the library, on the same
tasks and seeds, for \StudyCoreModels{} core models: \StudySEpisodes{} episodes and
\StudySFits{} fits on the development tasks.

\paragraph{Models and analysis.}
Tracks P and E run on \StudyModels{} open-weight instruction models from \StudyFamilies{}
families, 2B to 32B parameters, in BF16 at temperature 0.1, with a \StudyContext{}-token
context (16,384 for one model with a smaller native limit) and at most \StudyOutputCap{}
generated tokens per call. Every refusal, truncated reply and failed program stays in the
denominator. The position tasks cover only \StudyAlgorithms{} algorithms, and variants and
seeds are not independent of them: with three clusters a sign-flip test cannot go below
$p=0.25$. We therefore report P and S as effects per algorithm, without significance
claims. Track E weights its two sources equally.
